% Pursuing Stacks (cotes 134-2 to 134-8): archive cote 134-8, site page 14 (PDF page 15), typescript page 568.
% typescript. Draft compared: ps7.tex:647-707.
% AI-transcribed by Claude Opus 5.5 (claude-opus-5-5, adaptive thinking, effort high), Message Batches API, batch ps4 (msgbatch_012Aoc1NuskATPsyBeKSYK7J), 2026-11-12; not yet checked by eye.
\scanpage{134-8/14}
\mspage{568}

\xout{\unsure{it}} ty is a matter of course,  as far as the two last among the functors
(2),  together with the \inserted{two} internal operations,  are concerned,  as in
both contexts ( homology and \unsure{cbhomology|cohomology} formalism within (Cat) on the
one hand,  and the ``sweeping duality formalism'' on the other) \xout{these}
\xout{operations} the formulaire concerning these four operations
\[ f^{*} \ , \ \mathrm{R}f_{*} \ , \ \overset{L}{\otimes} \ , \ \mathrm{R}\underline{\mathrm{Hom}} \]
is no more, no less than just the relevant formulaire \xout{for} in the context
of \inserted{arbitrary} (commutatively) ringed topoi,  and maps between such.  The common
notation  $f_{!}$  ( occuring in $\mathrm{L}f_{!}$  in the (Cat) context,  in $\mathrm{R}f_{!}$ in the
``duality'' context)  is a very suggestive one,  for the least,  and I am
rather confident that \inserted[typed]{most}\xout{the} reflexes (concerning formal behaviour of $f_{!}$
with respect to the other operations)  coming from one context,  should
be OK too in the other.  Whether this is just a mere formal analogy,
or whether there is a deeper \inserted[typed]{relationship}\xout{\unsure{identity}}between the two kind of contexts,
I am at a loss at present to say.  It does'nt seem at all unlikely to
me that among arbitrary maps in (Cat),  one can single out some\hand{,} by
some kind of ``finite type'' condition,  for which a functor
\[ (*) \qquad f^{!} \ : \ D^{+}(B^{\wedge}_{k}) \longrightarrow D^{+}(A^{\wedge}_{k}) \]
\margin{\illegible{1}}
can be defined (where $f: A \longrightarrow B$),  right adjoint to the familiar $\mathrm{R}f_{*}$
functor ,  so that (2) can be completed to a sequence of \uline{four} functors
\[ (3) \qquad \mathrm{L}f_{!} \ , \ f^{*} \ , \ \mathrm{R}f_{*} \ , \ f^{!} \ , \]
forming a se\over{q}{q}uence of mutually \unsure{adjoiht|adjoint} functors between derived catego-%
ries,  in the usual sense (the functor immediately to the right of
another being it's right adjoint).  I faintly remember that Verdier
worked out such a formalism within the context of discrete or profinite
groups, or both,  in a Bourbaki talk\inserted{he gave}, this being inspired by the simi-%
lar work he did within the context of usual topological spaces.  In the
latter,  the finiteness condition required for a map \xout{f}
\[ f: X \longrightarrow Y \]
of topological spaces to give rise to a functor  $f^{!}$  (\xout{from} between\hand{,} say\hand{,}
the derived categories of the categories of abelian sheaves on Y and X)
is \inserted{mainly} that X should be locally embeddable in a product  $Y\mathrm{x}\underline{R}^{d}$  -  a rather
natural condition indeed!   As we are using (Cat) as a kind of algebrai
% NOTE: the last letter(s) of "algebrai..." run off the right edge of the scan
paradigm for the category of topological spaces,  this last example
for ``sweeping duality''  makes it rather plausible \deleted{indeed} that something
of the same kind should exist indeed in (Cat)  -  and likewise for the
context of \over{g}{g}roups  (which we may view as just particular cases of

\margin{14}
